\documentclass[answers,12pt,addpoints]{exam}

\usepackage{import}
\usepackage{amsmath,amssymb,amsthm,mathtools,bm}
\usepackage{tikz}
\usepackage[hidelinks]{hyperref}

% Optional local style; the document also compiles without this file.
\IfFileExists{C:/Users/prana/OneDrive/Desktop/MathNotes/style.tex}{%
  \import{C:/Users/prana/OneDrive/Desktop/MathNotes/}{style.tex}%
}{}

\providecommand{\name}{Pranav Tikkawar}
\providecommand{\course}{Statistics 652}
\providecommand{\assignment}{Homework 2}
\author{\name}
\title{\course\ -- \assignment}
\date{}

\providecommand{\E}{\mathbb{E}}
\providecommand{\R}{\mathbb{R}}
\providecommand{\Var}{\operatorname{Var}}
\providecommand{\IF}{\operatorname{IF}}
\providecommand{\GES}{\gamma^*}
\providecommand{\sgn}{\operatorname{sgn}}
\providecommand{\diag}{\operatorname{diag}}
\providecommand{\ind}{\mathbf{1}}
\providecommand{\ARE}{\operatorname{ARE}}
\renewcommand{\questionlabel}{2.\arabic{question}.}

\begin{document}
\maketitle

\section*{Notation and a preliminary lemma}
For a distribution function $F$ on $\R$, let
\[
Q_F(u)=\inf\{x:F(x)\geq u\},\qquad 0<u<1,
\]
and let $\Delta_x$ denote unit point mass at $x$. Define
\[
F_{\varepsilon,x}=(1-\varepsilon)F+\varepsilon\Delta_x,
\qquad 0\leq\varepsilon<1.
\]
The influence function and gross-error sensitivity are
\[
\IF(x;T,F)=\lim_{\varepsilon\downarrow0}
\frac{T(F_{\varepsilon,x})-T(F)}{\varepsilon},
\qquad
\GES(T,F)=\sup_{x\in\R}|\IF(x;T,F)|.
\]
All sampling statements concern independent identically distributed
observations unless otherwise stated. An asymptotic variance is the
variance of the root-$n$ limiting distribution, not that variance
divided by $n$.

\paragraph{Regular M-functional lemma.}
Let $t_0=T(F)$ solve
\[
\lambda_F(t_0)=0,\qquad \lambda_F(t)=\E_F\psi(X,t).
\]
Suppose $\lambda_F$ is differentiable at $t_0$, with
$A_F=\lambda_F'(t_0)\ne0$. For a fixed contamination point $x$, assume
$\psi(x,\cdot)$ is continuous at $t_0$ and that the selected roots
$t_\varepsilon=T(F_{\varepsilon,x})$ satisfy $t_\varepsilon\to t_0$.
Then
\[
\IF(x;T,F)=-A_F^{-1}\psi(x,t_0).
\]
\begin{proof}
The contaminated estimating equation is
\[
(1-\varepsilon)\lambda_F(t_\varepsilon)
+\varepsilon\psi(x,t_\varepsilon)=0.
\]
Put $d_\varepsilon=t_\varepsilon-t_0$. Differentiability gives
$\lambda_F(t_\varepsilon)=(A_F+r_\varepsilon)d_\varepsilon$, where
$r_\varepsilon\to0$. Hence
\[
\frac{d_\varepsilon}{\varepsilon}
=-\frac{\psi(x,t_\varepsilon)}
{(1-\varepsilon)(A_F+r_\varepsilon)}
\longrightarrow -A_F^{-1}\psi(x,t_0).
\]
\end{proof}
For a smooth population equation, a nonzero parameter derivative
and the implicit function theorem supply the required continuous
local root branch. For a vector parameter, the same argument uses a
nonsingular Jacobian. Differentiability of $\lambda_F$ does not by
itself license replacing $\lambda_F'$ by $\E_F\partial_t\psi$;
that interchange requires justification.

\newpage
\begin{questions}

% ============================================================
\question
Recall that the mean is $T(F)=\E_FX$ and the median is
$T(F)=F^{-1}(0.5)$.
\begin{parts}
\part For continuous $F$, express the trimmed and Winsorized means as
functionals $T(F)$.
\part For absolutely continuous $F$, derive the median's influence
function and gross-error sensitivity, and evaluate
$\E_F[\IF(X;T,F)]$ and $\E_F[\IF(X;T,F)^2]$.
\end{parts}

\begin{solution}
\paragraph{Assumptions.}
Use symmetric trimming with $0<\alpha<1/2$. For the median derivative,
assume $F$ is absolutely continuous and differentiable at its median
$m$, with $f(m)=F'(m)>0$.

\paragraph{Proof.}
\begin{parts}
\part
Let $a=Q_F(\alpha)$ and $b=Q_F(1-\alpha)$. Since $Q_F(U)$ has
CDF $F$ for $U\sim\operatorname{Uniform}(0,1)$, averaging the retained
central probability mass gives
\[
T_{\mathrm{tr},\alpha}(F)
=\frac{1}{1-2\alpha}\int_\alpha^{1-\alpha}Q_F(u)\,du.
\]
For continuous $F$, this equals
\[
\frac1{1-2\alpha}\int_a^b x\,dF(x).
\]
The Winsorizing map is $w_{a,b}(x)=\min\{b,\max\{a,x\}\}$.
Its expectation is
\[
T_{\mathrm W,\alpha}(F)
=\alpha a+\int_\alpha^{1-\alpha}Q_F(u)\,du+\alpha b.
\]
No renormalization is needed because Winsorizing moves the tail
mass to the cutoffs rather than deleting it. These quantile
representations also define the functionals for distributions with
atoms and require no tail moment assumptions.

\part
Write $m_\varepsilon=Q_{F_{\varepsilon,x}}(1/2)$.
Because $F'(m)>0$, for all sufficiently small $\delta>0$,
\[
F(m-\delta)<1/2<F(m+\delta).
\]
The same strict inequalities hold under sufficiently small
contamination. Therefore $m_\varepsilon\to m$.
Fix $x\ne m$. The contaminating atom then stays on one side of a
neighborhood of $m$, so
\[
m_\varepsilon=Q_F(u_\varepsilon),\qquad
u_\varepsilon
=\frac{\frac12-\varepsilon\ind\{x<m\}}{1-\varepsilon}.
\]
In particular,
\[
u_\varepsilon-\frac12
=\varepsilon\left(\frac12-\ind\{x<m\}\right)+O(\varepsilon^2).
\]
To differentiate the quantile, put $y=Q_F(u)$ as $u\to1/2$.
The strict inequalities around $m$ imply $y\to m$, and continuity
of $F$ gives $F(y)=u$. Consequently,
\[
\frac{Q_F(u)-m}{u-1/2}
=\frac{y-m}{F(y)-F(m)}\longrightarrow\frac1{f(m)}.
\]
Factoring the contamination difference quotient now proves
\[
\IF(x;T,F)=\frac{\frac12-\ind\{x<m\}}{f(m)}
=\frac{\sgn(x-m)}{2f(m)},\qquad x\ne m.
\]
At $x=m$, the contaminated median remains $m$, giving influence zero.
The median influence is therefore bounded by $1/(2f(m))$.
Since $\Pr(X<m)=\Pr(X>m)=1/2$, its positive and negative values
have equal probabilities, and its square is constant almost surely.
\end{parts}

\paragraph{Final answer.}
\[
\boxed{T_{\mathrm{tr},\alpha}(F)
=\frac{\int_\alpha^{1-\alpha}Q_F(u)\,du}{1-2\alpha}}
\]
\[
\boxed{T_{\mathrm W,\alpha}(F)
=\alpha Q_F(\alpha)+\int_\alpha^{1-\alpha}Q_F(u)\,du
+\alpha Q_F(1-\alpha)}
\]
\[
\boxed{\IF(x;\operatorname{Med},F)
=\frac{\sgn(x-m)}{2f(m)},\quad \sgn(0)=0;
\qquad \GES=\frac1{2f(m)}}
\]
\[
\boxed{\E_F\IF(X;T,F)=0,\qquad
\E_F[\IF(X;T,F)^2]=\frac1{4f(m)^2}.}
\]
\end{solution}

% ============================================================
\question
\begin{parts}
\part Show that the mode is not qualitatively robust.
\part Show that the trimmed and Winsorized means can be expressed as
L-statistics: find $a_1,\ldots,a_n$ such that
$T_n=\sum_{i=1}^n a_i x_{(i)}$.
\end{parts}

\begin{solution}
\paragraph{Assumptions.}
For (a), interpret a mode estimator as a sequence consistently
estimating the unique population mode at the distributions used
below. Qualitative robustness refers to weak perturbations of the
sampling distribution. For (b), use integer trimming of $k$
observations from each tail, with $0\leq2k<n$.

\paragraph{Proof.}
\begin{parts}
\part
Let $F=N(0,1)$ with density $\phi$ and mode zero. Fix $a\ne0$, and put
\[
K(u)=\frac34(1-u^2)\ind\{|u|\leq1\},\qquad h_\varepsilon=\varepsilon^2,
\]
\[
g_\varepsilon(t)=(1-\varepsilon)\phi(t)
+\frac{\varepsilon}{h_\varepsilon}
K\left(\frac{t-a}{h_\varepsilon}\right).
\]
For its distribution $G_\varepsilon$,
\[
d_{\mathrm{TV}}(G_\varepsilon,F)\leq\varepsilon,
\quad\text{so}\quad G_\varepsilon\Rightarrow F.
\]
Nevertheless, $g_\varepsilon(a)=(1-\varepsilon)\phi(a)+3/(4\varepsilon)$
tends to infinity, whereas outside $[a-h_\varepsilon,a+h_\varepsilon]$
the density is bounded by $\sup\phi$. All global modes lie inside
that interval for small $\varepsilon$. There,
\[
g_\varepsilon''(t)
=(1-\varepsilon)\phi''(t)-\frac32\varepsilon^{-5}<0
\]
for sufficiently small $\varepsilon$. Hence the mode is unique and
\[
m_\varepsilon=T_{\mathrm{mode}}(G_\varepsilon)\to a\ne0.
\]
Thus the mode functional is not weakly continuous.

To establish non-robustness of an estimator sequence, let $T_n$ be
consistent at $F$ and every fixed $G_\varepsilon$ above. Write $d_\pi$
for Prokhorov distance. Qualitative robustness at $F$ would imply:
for each $\zeta>0$ there is $\delta>0$ such that
\[
d_\pi(G,F)<\delta\ \Longrightarrow\
d_\pi(\mathcal L_G(T_n),\mathcal L_F(T_n))<\zeta
\]
uniformly in $n$ (the argument also applies to the asymptotic version).
Choose $0<\zeta<\min\{|a|,1\}/2$, then fix small $\varepsilon<\delta$
with $\min\{|m_\varepsilon|,1\}>\zeta$. Since
$d_\pi(G_\varepsilon,F)\leq\varepsilon$, robustness gives the above
bound. Consistency and passage to $n\to\infty$ imply
\[
d_\pi(\Delta_{m_\varepsilon},\Delta_0)\leq\zeta.
\]
But $d_\pi(\Delta_s,\Delta_t)=\min\{|s-t|,1\}$, a contradiction.

Such consistent estimators exist for these distributions. For example,
use a Gaussian-kernel density estimator with bandwidth $b_n=n^{-1/4}$
and a measurable global maximizer as the mode estimator. Integration
by parts and the empirical-CDF bound give
\[
\sup_t|\widehat f_n(t)-\E\widehat f_n(t)|
\leq \frac{\|\phi'\|_1}{b_n}\sup_y|F_n(y)-F(y)|
=O_p(n^{-1/4}).
\]
For each fixed density above, bounded uniform continuity gives
uniformly vanishing bias. Each density tends to zero at infinity and
has a unique mode; hence that maximum is separated from the supremum
outside any fixed neighborhood. Uniform density consistency implies
mode consistency. Therefore the failure is one of qualitative
robustness of consistent mode estimation, not merely a discontinuity
statement without an estimator interpretation.

\part
For the trimmed mean,
\[
T_{\mathrm{tr},n}=\frac1{n-2k}\sum_{i=k+1}^{n-k}x_{(i)}.
\]
For the Winsorized mean,
\[
T_{\mathrm W,n}=\frac1n\left[
kx_{(k+1)}+\sum_{i=k+1}^{n-k}x_{(i)}+kx_{(n-k)}\right].
\]
Collecting the coefficient of each order statistic gives the desired
L-statistic representations. The formula below also handles the
case $k+1=n-k$.
\end{parts}

\paragraph{Final answer.}
The mode functional is not weakly continuous, and any mode estimator
consistent at the displayed distributions is not qualitatively robust
at $N(0,1)$. The L-statistic coefficients are
\[
\boxed{a_i^{\mathrm{tr}}=
\frac{\ind\{k+1\leq i\leq n-k\}}{n-2k}}
\]
and
\[
\boxed{a_i^{\mathrm W}=\frac1n\left[
\ind\{k+1\leq i\leq n-k\}
+k\ind\{i=k+1\}+k\ind\{i=n-k\}\right].}
\]
For an exact empirical quantile-integral definition rather than
integer trimming, let $I_i=((i-1)/n,i/n]$ and let $\lambda$ be length.
Then the corresponding coefficients are
\[
a_i^{\mathrm{tr},\alpha}
=\frac{\lambda(I_i\cap[\alpha,1-\alpha])}{1-2\alpha},
\]
\[
a_i^{\mathrm W,\alpha}
=\lambda(I_i\cap[\alpha,1-\alpha])
+\alpha\ind\{i=\lceil n\alpha\rceil\}
+\alpha\ind\{i=\lceil n(1-\alpha)\rceil\}.
\]
These exact-quantile conventions need not coincide with integer
trimming at boundary ranks.
\end{solution}

% ============================================================
\question
For a location M-estimate, show (i) $\rho(x,t)=\rho(x-t)$ with
$\rho(-r)=\rho(r)$, and (ii) $\psi(x,t)=\psi(x-t)$ with
$\psi(-r)=-\psi(r)$.

\begin{solution}
\paragraph{Assumptions.}
Use a translation-invariant, reflection-symmetric loss:
\[
\rho(x+a,t+a)=\rho(x,t),\qquad \rho(-x,-t)=\rho(x,t).
\]
Assume differentiability when an ordinary score is used. Reflection
symmetry is an additional assumption: it does not follow from
location equivariance alone. For example,
$\rho_\tau(r)=r(\tau-\ind\{r<0\})$ is a location loss but is not
even for $\tau\ne1/2$.

\paragraph{Proof.}
Taking $a=-t$ gives $\rho(x,t)=\rho(x-t,0)$. Define
$\rho_0(r)=\rho(r,0)$. Then
\[
\rho(x,t)=\rho_0(x-t),\qquad
\rho_0(-r)=\rho(-r,0)=\rho(r,0)=\rho_0(r).
\]
Let $\psi_0=\rho_0'$. Differentiating the evenness identity yields
$-\psi_0(-r)=\psi_0(r)$, so $\psi_0$ is odd and $\psi_0(0)=0$.
The minimizing first-order equation is
\[
\sum_i\psi_0(x_i-t)=0.
\]
If $\psi(x,t)$ is defined as $\partial_t\rho(x,t)$, it equals
$-\psi_0(x-t)$; this convention gives the same roots and oddness.
For convex nonsmooth losses, replace the derivative equation by
$0\in\sum_i\partial\rho_0(x_i-t)$.

If $F$ is symmetric about $\theta$ and the score is integrable,
oddness also implies $\E_F\psi_0(X-\theta)=0$. Consequently the
functional is Fisher consistent whenever that is the unique selected
root.

\paragraph{Final answer.}
Under translation invariance and reflection symmetry,
\[
\boxed{\rho(x,t)=\rho_0(x-t),\quad \rho_0(-r)=\rho_0(r),
\qquad \psi_0(-r)=-\psi_0(r).}
\]
\end{solution}

% ============================================================
\question
Derive the asymptotic relative efficiency of the mean to the median
at (i) $X\sim\mu+\sigma t_\nu$ and (ii)
$X\sim(1-\eta)N(\mu,\sigma^2)+\eta N(\mu,9\sigma^2)$.

\begin{solution}
\paragraph{Assumptions and convention.}
Define $\ARE(A,B)=V_B/V_A$, where $V$ is root-$n$ asymptotic
variance. Thus an ARE exceeding one favors $A$.
For a common mean and unique median $\mu$, finite variance $v$,
and differentiability at $\mu$ with $f(\mu)>0$, the CLT and 2.9 give
\[
V_{\mathrm{mean}}=v,\qquad V_{\mathrm{median}}=\frac1{4f(\mu)^2},
\qquad \ARE(\bar X,\widehat m)=\frac1{4vf(\mu)^2}.
\]

\paragraph{Proof.}
\begin{parts}
\part
For $\nu>2$,
\[
v=\sigma^2\frac\nu{\nu-2},\qquad f_X(\mu)=\frac{c_\nu}\sigma,
\qquad
c_\nu=\frac{\Gamma((\nu+1)/2)}{\sqrt{\nu\pi}\Gamma(\nu/2)}.
\]
Substitution gives
\[
\ARE(\bar X,\widehat m)
=\frac{\nu-2}{4\nu c_\nu^2}
=\frac{\pi(\nu-2)}4
\left[\frac{\Gamma(\nu/2)}{\Gamma((\nu+1)/2)}\right]^2.
\]
The gamma-ratio asymptotic
$\Gamma(z+1/2)/\Gamma(z)\sim z^{1/2}$ gives the limit $\pi/2$.
For example, the ARE is $\pi^2/16$ at $\nu=3$, $8/9$ at $\nu=4$,
and approximately $1.04093$ at $\nu=5$. The crossover is
$\nu\approx4.67879$.

For $0<\nu\leq2$, the ordinary ratio of finite root-$n$ variances
is not defined for the sample mean. The median remains root-$n$
normal for every $\nu>0$. Student tails are regularly varying with
index $-\nu$: for $0<\nu<2$, the generalized CLT gives
\[
n^{1-1/\nu}(\bar X-\mu)\xrightarrow{d}S_\nu,
\]
where $S_\nu$ is a nondegenerate symmetric stable law with the
appropriate scale. Thus the mean is slower than root-$n$ for
$1<\nu<2$, does not concentrate for $\nu=1$, and has increasing
spread for $0<\nu<1$. At $\nu=2$, its concentration rate is of order
$\sqrt{n/\log n}$, again slower than root-$n$.
Accordingly its efficiency relative to the median is zero under
the extended rate-based convention. For $\nu\leq1$, $\mu$ denotes
location, not a finite expectation.

\part
The component means agree, so
\[
v=(1-\eta)\sigma^2+\eta(9\sigma^2)=\sigma^2(1+8\eta).
\]
Symmetry gives median $\mu$, and
\[
f_X(\mu)=\frac{1-\eta}{\sigma\sqrt{2\pi}}
+\frac\eta{3\sigma\sqrt{2\pi}}
=\frac{1-\frac23\eta}{\sigma\sqrt{2\pi}}.
\]
Substitution yields the expression below. Its denominator apart from
the factor two is $D(\eta)=(1+8\eta)(1-\frac23\eta)^2$, with
\[
D'(\eta)=\frac43(1-\tfrac23\eta)(5-12\eta).
\]
Thus the ARE decreases until $\eta=5/12$ and then increases; it is
not monotone on $[0,1]$. Solving $\ARE=1$ gives approximately
$0.100592$ and $0.814126$. At $\eta=0$ and $1$, the distribution
is a single normal and the ARE is $\pi/2$.
\end{parts}

\paragraph{Final answer.}
\[
\boxed{\ARE(\bar X,\widehat m)
=\frac{\pi(\nu-2)}4
\left[\frac{\Gamma(\nu/2)}{\Gamma((\nu+1)/2)}\right]^2,
\quad\nu>2.}
\]
For $\nu\leq2$, finite root-$n$ variance ARE is unavailable; the
extended rate-based efficiency of the mean relative to the median
is zero.
\[
\boxed{\ARE(\bar X,\widehat m)
=\frac\pi{2(1+8\eta)(1-\frac23\eta)^2},\quad0\leq\eta\leq1.}
\]
The median is more efficient for
$0.100592\ldots<\eta<0.814126\ldots$.
\end{solution}

% ============================================================
\question
Let $\widehat\theta_n$ minimize $\sum_i\rho(x_i;\theta)$, with
$\rho$ the negative log-likelihood of a regular one-parameter family.
Show directly, without MLE asymptotic theory, that at $F=F_\theta$,
$\E_F[\IF(X;T,F)^2]=1/I(\theta;X)$.

\begin{solution}
\paragraph{Assumptions.}
Assume an identifiable regular family with common support,
differentiation under the integral justified up to second order,
and $0<I(\theta;X)<\infty$. Assume the population score equation
is continuously differentiable near the model parameter and under
small contamination. The functional follows the local minimizer
branch through $\theta$; if defined by global minimization, assume
its minimizer stays on this branch. Expected log-likelihood
differences below are assumed well-defined.

\paragraph{Proof.}
Put $s_t(x)=\partial_t\log f_t(x)$ and
$\psi(x,t)=\partial_t\rho(x;t)=-s_t(x)$. Fisher consistency follows
from
\[
\E_\theta\rho(X;t)-\E_\theta\rho(X;\theta)
=D_{\mathrm{KL}}(F_\theta\Vert F_t)\geq0,
\]
with equality only at $t=\theta$ by identifiability.
Thus $T(F_\theta)=\theta$.
Differentiating $\int f_\theta=1$ once and twice gives
\[
\E_\theta s_\theta(X)=0,
\qquad
\E_\theta[\partial_\theta s_\theta(X)+s_\theta(X)^2]=0.
\]
Consequently,
\[
A_{F_\theta}
=-\E_\theta\partial_\theta s_\theta(X)
=\E_\theta s_\theta(X)^2=I(\theta;X)>0.
\]
The implicit function theorem supplies a differentiable local root
branch, and the preliminary lemma gives
\[
\IF(x;T,F_\theta)=I(\theta;X)^{-1}s_\theta(x).
\]
Squaring and integrating proves the result directly.

\paragraph{Final answer.}
\[
\boxed{\IF(x;T,F_\theta)=\frac{s_\theta(x)}{I(\theta;X)},
\qquad
\E_\theta[\IF(X;T,F_\theta)^2]=\frac1{I(\theta;X)}.}
\]
For a vector parameter and nonsingular information matrix, the
corresponding identity is
$\E_\theta[\IF\IF^\top]=I(\theta)^{-1}$.
\end{solution}

% ============================================================
\question
For a distribution with a finite second moment, express its variance
$\sigma_F^2$ as a functional, find its influence function and
gross-error sensitivity, sketch the influence function, evaluate its
first two moments, and derive the sample variance's asymptotic
variance by the CLT.

\begin{solution}
\paragraph{Assumptions and notation.}
Let $\mu=\E_FX$, $v=\sigma_F^2$, and
$\mu_4=\E_F[(X-\mu)^4]$. The functional is defined on
$\mathcal F_2=\{F:\int x^2\,dF(x)<\infty\}$.
For the CLT assume $\mu_4<\infty$.
Here $\sigma_F^2$ is the variance; the wording ``standard deviation''
and centering by $\sigma$ in 2.6 are interpreted consistently with
the sample variance as $\sigma_F^2$.

\paragraph{Proof.}
The functional is
\[
T(F)=\int x^2\,dF(x)-\left(\int x\,dF(x)\right)^2.
\]
The contaminated mean is $\mu+\varepsilon(x-\mu)$, and the mixture
variance identity gives exactly
\[
T(F_{\varepsilon,x})=(1-\varepsilon)v
+\varepsilon(1-\varepsilon)(x-\mu)^2.
\]
Differentiating yields
$\IF(x;T,F)=(x-\mu)^2-v$. This diverges as $|x|\to\infty$.
Its graph is an upward-opening parabola with vertex $(\mu,-v)$ and,
for $v>0$, zeros $\mu\pm\sqrt v$.
\begin{center}
\begin{tikzpicture}[x=1.35cm,y=0.7cm]
\draw[->] (-2.2,0)--(2.4,0) node[right] {$z=(x-\mu)/\sqrt v$};
\draw[->] (0,-1.5)--(0,3.6) node[above] {$\IF(x;T,F)/v$};
\draw[thick,domain=-2:2,samples=70] plot (\x,{\x*\x-1});
\node[below left] at (0,-1) {$-1$};
\node[below] at (-1,0) {$-1$};
\node[below] at (1,0) {$1$};
\end{tikzpicture}
\end{center}
The displayed standardized sketch assumes $v>0$; when $v=0$, the
unscaled graph is $(x-\mu)^2$.
Taking moments,
\[
\E_F\IF(X;T,F)=\E_F[(X-\mu)^2]-v=0,
\]
\[
\E_F[\IF(X;T,F)^2]
=\E_F[((X-\mu)^2-v)^2]=\mu_4-v^2.
\]
The last quantity is infinite if $\mu_4=\infty$.

For $s_n^2=n^{-1}\sum_i(X_i-\bar X)^2$, the exact identity
\[
s_n^2=\frac1n\sum_i(X_i-\mu)^2-(\bar X-\mu)^2
\]
implies
\[
\sqrt n(s_n^2-v)
=\frac1{\sqrt n}\sum_i[(X_i-\mu)^2-v]
-\sqrt n(\bar X-\mu)^2.
\]
By the mean CLT, $\bar X-\mu=O_p(n^{-1/2})$, so the final term is
$o_p(1)$. The CLT for $(X_i-\mu)^2$ and Slutsky's theorem give
\[
\sqrt n(s_n^2-v)\xrightarrow{d}N(0,\mu_4-v^2).
\]
A zero limiting variance is allowed as a degenerate normal limit.

\paragraph{Final answer.}
\[
\boxed{T(F)=\E_FX^2-(\E_FX)^2,\quad
\IF(x;T,F)=(x-\mu)^2-v,\quad \GES=\infty.}
\]
\[
\boxed{\E_F\IF=0,\qquad \E_F\IF^2=\mu_4-v^2
=\operatorname{AVar}(s_n^2).}
\]
For normal $F$, this asymptotic variance is $2\sigma^4$.
\end{solution}

% ============================================================
\question
Express the standard deviation $\sigma$ as a functional, find and
sketch its influence function, and determine whether the sample
standard deviation's asymptotic variance equals $\E_F[\IF^2]$.

\begin{solution}
\paragraph{Assumptions.}
Let $F\in\mathcal F_2$ and $\sigma>0$. For asymptotic normality,
assume $\mu_4=\E_F[(X-\mu)^4]<\infty$.

\paragraph{Proof.}
The standard deviation functional is the square root of the variance:
\[
T(F)=\left[\int x^2\,dF(x)-\left(\int x\,dF(x)\right)^2\right]^{1/2}.
\]
Applying the ordinary chain rule to the contamination derivative
from 2.6 gives
\[
\IF(x;T,F)=\frac{(x-\mu)^2-\sigma^2}{2\sigma}.
\]
This is an upward-opening parabola with vertex $(\mu,-\sigma/2)$
and zeros $\mu\pm\sigma$; it has infinite gross-error sensitivity.
\begin{center}
\begin{tikzpicture}[x=1.35cm,y=1cm]
\draw[->] (-2.2,0)--(2.4,0) node[right] {$z=(x-\mu)/\sigma$};
\draw[->] (0,-0.9)--(0,2.2) node[above] {$\IF(x;T,F)/\sigma$};
\draw[thick,domain=-2:2,samples=70] plot (\x,{(\x*\x-1)/2});
\node[below left] at (0,-0.5) {$-1/2$};
\node[below] at (-1,0) {$-1$};
\node[below] at (1,0) {$1$};
\end{tikzpicture}
\end{center}
Direct calculation gives
\[
\E_F\IF=0,\qquad
\E_F\IF^2=\frac{\mu_4-\sigma^4}{4\sigma^2}.
\]
Applying the delta method to 2.6 with $g(v)=\sqrt v$ yields
\[
\sqrt n(s_n-\sigma)\xrightarrow{d}
N\left(0,\frac{\mu_4-\sigma^4}{4\sigma^2}\right).
\]
At $\sigma=0$, $F=\Delta_\mu$, and exactly
\[
T(F_{\varepsilon,x})=\sqrt{\varepsilon(1-\varepsilon)}\,|x-\mu|.
\]
For $x\ne\mu$, the contamination difference quotient is infinite
in the limit, so there is no finite regular influence function there.

\paragraph{Final answer.}
For $\sigma>0$,
\[
\boxed{\IF(x;T,F)=\frac{(x-\mu)^2-\sigma^2}{2\sigma},\qquad
\operatorname{AVar}(s_n)=\E_F\IF^2
=\frac{\mu_4-\sigma^4}{4\sigma^2}.}
\]
For normal $F$, the asymptotic variance is $\sigma^2/2$.
\end{solution}

% ============================================================
\question
Derive the density of $X=\mu+\sigma RZ$, where $\mu\in\R$,
$\sigma>0$, $Z\sim N(0,1)$, $R>0$ has density $f_R$, and $R,Z$
are independent.

\begin{solution}
\paragraph{Assumptions.}
Only positivity of $R$, the stated independence, and existence of
the mixing density are needed; no moments of $R$ are required.

\paragraph{Proof.}
Conditioning gives $X\mid R=r\sim N(\mu,\sigma^2r^2)$ and
\[
f_{X\mid R}(x\mid r)=\frac1{\sigma r\sqrt{2\pi}}
\exp\left[-\frac{(x-\mu)^2}{2\sigma^2r^2}\right].
\]
Integrating this conditional density against $f_R$ gives the result.
Tonelli's theorem verifies normalization:
\[
\int_\R f_X(x)\,dx
=\int_0^\infty f_R(r)
\left[\int_\R f_{X\mid R}(x\mid r)\,dx\right]dr=1.
\]
The density is symmetric about $\mu$. Its mixture version at the
center has value $\E[R^{-1}]/(\sigma\sqrt{2\pi})$, possibly infinite.

\paragraph{Final answer.}
\[
\boxed{f_X(x)=\int_0^\infty
\frac{\exp[-(x-\mu)^2/(2\sigma^2r^2)]}{\sigma r\sqrt{2\pi}}
f_R(r)\,dr.}
\]
For a general mixing law $H_R$, replace $f_R(r)\,dr$ by $dH_R(r)$;
this includes discrete scale mixtures.
\end{solution}

% ============================================================
\question
(Optional.) Prove consistency and asymptotic normality of the sample
median: if $\mu$ is the population median and $f(\mu)>0$, establish
the limiting distribution of $\sqrt n(\widehat\mu_n-\mu)$.

\begin{solution}
\paragraph{Assumptions.}
Assume $F$ is continuous near its unique median $\mu$ and is
differentiable at $\mu$ with derivative $f(\mu)>0$.
Let $k_n=\lceil n/2\rceil$ and initially set
$\widehat\mu_n=X_{(k_n)}$. The centered limiting normal has mean
zero; the mean $\mu$ printed in the assignment's limiting law is
corrected accordingly.

\paragraph{Proof of consistency.}
For every $\delta>0$,
\[
F(\mu-\delta)<1/2<F(\mu+\delta).
\]
If $\widehat\mu_n>\mu+\delta$, then
$F_n(\mu+\delta)<k_n/n$. By the law of large numbers and
$k_n/n\to1/2$, this event has probability tending to zero.
Likewise $\widehat\mu_n<\mu-\delta$ implies
$F_n(\mu-\delta)\geq k_n/n$, again with probability tending to zero.
Therefore $\widehat\mu_n\xrightarrow{p}\mu$.

\paragraph{Proof of asymptotic normality.}
Fix $t\in\R$ and set
\[
y_n=\mu+t/\sqrt n,\qquad
p_n=F(y_n)=\frac12+\frac{f(\mu)t}{\sqrt n}+o(n^{-1/2}).
\]
The count $B_n=\sum_i\ind\{X_i\leq y_n\}$ is
$\operatorname{Bin}(n,p_n)$. Since $p_n\to1/2$, the standardized
Bernoulli triangular array satisfies Lindeberg's condition, giving
\[
Z_n=\frac{B_n-np_n}{\sqrt{np_n(1-p_n)}}\xrightarrow{d}N(0,1).
\]
The order-statistic event satisfies
\[
\Pr(\widehat\mu_n\leq y_n)
=\Pr(B_n\geq k_n)
=1-\Pr\left(Z_n\leq
\frac{k_n-1-np_n}{\sqrt{np_n(1-p_n)}}\right).
\]
The threshold converges to $-2f(\mu)t$. Convergence to the continuous
normal CDF permits this moving threshold, so the last expression
tends to $\Phi(2f(\mu)t)$.

For even $n=2k$, the same argument applies to ranks $k$ and $k+1$.
For $M_n=(X_{(k)}+X_{(k+1)})/2$,
\[
\Pr(\sqrt n(X_{(k+1)}-\mu)\leq t)
\leq\Pr(\sqrt n(M_n-\mu)\leq t)
\leq\Pr(\sqrt n(X_{(k)}-\mu)\leq t).
\]
Both outer terms have the same limit, proving the result for the
usual averaged median without an order-statistic spacing assumption.

\paragraph{Final answer.}
\[
\boxed{\widehat\mu_n\xrightarrow{p}\mu,\qquad
\sqrt n(\widehat\mu_n-\mu)\xrightarrow{d}
N\left(0,\frac1{4f(\mu)^2}\right).}
\]
\end{solution}

% ============================================================
\question
Show that if $X_1,X_2$ are independent
$\operatorname{Cauchy}(\mu,\sigma^2)$ variables, then
$(X_1+X_2)/2$ has the same distribution.

\begin{solution}
\paragraph{Assumptions and parameterization.}
Let $s>0$ denote the Cauchy scale under the assignment's convention.
Its density is $s/\{\pi[(x-\mu)^2+s^2]\}$. The second parameter
in the problem is a scale convention, not a finite variance.

\paragraph{Proof.}
The characteristic function is $\phi_X(t)=\exp(i\mu t-s|t|)$.
Independence gives
\[
\phi_{(X_1+X_2)/2}(t)
=\phi_{X_1}(t/2)\phi_{X_2}(t/2)
=\exp(i\mu t-s|t|)=\phi_X(t).
\]
Uniqueness of characteristic functions proves the claim.

\paragraph{Final answer.}
\[
\boxed{(X_1+X_2)/2\overset d=X_1
\sim\operatorname{Cauchy}(\mu,\sigma^2).}
\]
The same calculation gives $\bar X_n\overset d=X_1$ for every $n$.
\end{solution}

% ============================================================
\question
A standardized symmetric stable variable with index $0<\alpha\leq2$
has characteristic function $\phi_X(t)=\exp(-|t|^\alpha)$.
For i.i.d. copies $X_i$ and $Y_i=\mu+\sigma X_i$, find the
characteristic function of $\bar Y_n=n^{-1}\sum_iY_i$ and describe
how its spread depends on $n$.

\begin{solution}
\paragraph{Assumptions.}
Let $\sigma>0$ and use the stated characteristic-function
standardization. Spread is measured by scale or fixed interquantile
ranges, not by variance when it is infinite.

\paragraph{Proof.}
Independence gives
\begin{align*}
\phi_{\bar Y_n}(t)
&=e^{i\mu t}\prod_{i=1}^n
\exp\left[-\left|\frac{\sigma t}{n}\right|^\alpha\right]\\
&=\exp\left[i\mu t-\sigma^\alpha n^{1-\alpha}|t|^\alpha\right].
\end{align*}
This is the characteristic function of
$\mu+\sigma n^{1/\alpha-1}X$. Therefore its interquantile ranges
are multiplied by $n^{1/\alpha-1}$.

\paragraph{Final answer.}
\[
\boxed{\phi_{\bar Y_n}(t)
=\exp\left[i\mu t-\sigma^\alpha n^{1-\alpha}|t|^\alpha\right],
\qquad
\bar Y_n\overset d=\mu+\sigma n^{1/\alpha-1}X.}
\]
Spread decreases for $\alpha>1$, stays constant for $\alpha=1$,
and increases for $0<\alpha<1$.
At $\alpha=2$, the standard law is $N(0,2)$ and the scale decreases
as $n^{-1/2}$. For $\alpha<2$, the variance is infinite; for
$\alpha\leq1$, $\mu$ is location rather than a finite mean.
\end{solution}

% ============================================================
\question
Prove that $\E_F[\IF(X;T,F)]=0$ whenever $T$ is an M-functional.

\begin{solution}
\paragraph{Assumptions.}
Let $T(F)=t_0$ solve $\lambda_F(t_0)=\E_F\psi(X,t_0)=0$, with
$\E_F|\psi(X,t_0)|<\infty$. Assume
$A_F=\lambda_F'(t_0)\ne0$, continuous selected contaminated roots,
and continuity of $\psi(x,\cdot)$ at $t_0$ for $F$-almost every $x$.
These are regularity conditions for the influence function to exist;
singular or nonunique nonregular functionals need not possess one.

\paragraph{Proof.}
The preliminary lemma gives, for $F$-almost every $x$,
\[
\IF(x;T,F)=-A_F^{-1}\psi(x,t_0).
\]
The integrability assumption makes this influence function integrable.
Taking expectations and using the population estimating equation,
\[
\E_F\IF(X;T,F)
=-A_F^{-1}\E_F\psi(X,t_0)=0.
\]
The argument only requires differentiability of the population
expectation. For example, with the median score
$\psi(x,t)=\sgn(x-t)$ and continuous $F$,
\[
\lambda_F(t)=1-2F(t),\qquad A_F=-2f(m).
\]
The pointwise score derivative is zero off its jump, so taking its
expectation would incorrectly miss this nonzero population derivative.
For Huber's Lipschitz score, bounded difference quotients instead
justify dominated differentiation, yielding
\[
A_F=-\Pr_F(|X-t_0|<c)
\]
when there is no mass at the two kinks.
The vector case follows identically with a nonsingular Jacobian.

\paragraph{Final answer.}
For every regular M-functional satisfying the stated conditions,
\[
\boxed{\E_F[\IF(X;T,F)]=0.}
\]
\end{solution}

% ============================================================
\question
Show that Tukey's M-estimate of location does not correspond to a
maximum likelihood estimate under a symmetric location model.

\begin{solution}
\paragraph{Assumptions and criterion.}
For $c>0$, use Tukey's bisquare loss
\[
\rho_c(r)=\begin{cases}
\dfrac{c^2}{6}\left[1-\left(1-\dfrac{r^2}{c^2}\right)^3\right],
&|r|\leq c,\\[1ex]
\dfrac{c^2}{6},&|r|>c,
\end{cases}
\]
with score $\psi_c(r)=r(1-r^2/c^2)^2\ind\{|r|\leq c\}$.
A likelihood representation of the defining loss or score is
considered first. To rule out even an alternative regular likelihood
having the same estimator for all samples, assume its symmetric
location density $g$ is positive and continuously differentiable.

\paragraph{Proof.}
If the loss were a negative log-density up to a positive scale and
an additive constant, then
\[
g(r)=C\exp[-a\rho_c(r)],\qquad C>0,\ a>0.
\]
For $|r|>c$, this is the positive constant $Ce^{-ac^2/6}$.
Hence it cannot integrate to one. The same conclusion follows from
matching likelihood scores: $-(\log g)'=a\psi_c$ integrates to the
same improper expression. This obstruction to a loss/score likelihood
representation does not require smoothness assumptions on a candidate
density beyond the proposed representation itself.

There is also a direct estimator-level contradiction for any positive
continuously differentiable symmetric density, even without requiring
proportional scores. Consider the three observations $(0,0,r)$ with
$r>2c$, and write $C_c=c^2/6$. The Tukey objective is
\[
L(t)=2\rho_c(t)+\rho_c(r-t).
\]
If $|t|<c$, then $|r-t|>c$, so $L(t)=2\rho_c(t)+C_c$, uniquely
minimized there at $t=0$. If $|t|\geq c$, then
$L(t)\geq2C_c>C_c=L(0)$. Thus its unique global minimizer is zero.

If this estimator were also the MLE under $g$ for every sample,
zero would maximize $g(-t)^2g(r-t)$. The derivative of its
log-likelihood at zero must vanish:
\[
-2\frac{g'(0)}{g(0)}-\frac{g'(r)}{g(r)}=0.
\]
Symmetry and differentiability imply $g'(0)=0$, hence $g'(r)=0$
for every $r>2c$. Thus $g$ is a positive constant on $(2c,\infty)$,
contradicting integrability.

\paragraph{Final answer.}
\[
\boxed{\text{Tukey's bisquare loss/score has no proper
location likelihood representation.}}
\]
Moreover, no positive continuously differentiable symmetric location
density has the same location MLE as Tukey's global M-estimator for
all samples. The bounded loss, equivalently the score's exact tail
cutoff, is the obstruction. Huber's linearly growing tail loss does
not have this obstruction: its negative exponential is integrable.
\end{solution}

% ============================================================
\question
For the location-scale double-exponential distribution, find the MLE,
show that the scale estimate has an unbounded influence function,
and find the finite-sample breakdown point of the scale estimate.

\begin{solution}
\paragraph{Assumptions and parameterization.}
Use
\[
f(x;\mu,b)=\frac1{2b}\exp\left(-\frac{|x-\mu|}{b}\right),\quad b>0.
\]
Its standard deviation is $\sqrt2 b$. For influence calculations,
assume a finite first moment, absolute continuity, and $f(m)>0$
at the median. For breakdown, take $n\geq2$ and an original sample
with positive fitted scale. For even $n$, select the averaged median.

\paragraph{Proof of the MLE.}
The log-likelihood is
\[
\ell(\mu,b)=-n\log(2b)-\frac1b\sum_i|x_i-\mu|.
\]
For fixed $\mu$, its maximizing scale is
$\widehat b(\mu)=n^{-1}\sum_i|x_i-\mu|$ when the sum is positive.
Profiling reduces the location problem to minimizing the sum of
absolute deviations, whose minimizers are the sample medians.
The minimum sum is the same for every median in the even-sample
median interval. If all data coincide, the likelihood diverges as
$b\downarrow0$ and no positive-scale MLE exists; extend the fitted
scale continuously by zero for breakdown calculations.

\paragraph{Proof of the influence function.}
Let $H(t)=\E_F|X-t|$, $m=m(F)$, and $B(F)=H(m)$.
The bound
\[
\left||X-t-h|-|X-t|\right|\leq|h|
\]
justifies dominated differentiation at points without an atom, so
$H'(t)=2F(t)-1$ and $H'(m)=0$.
Problem 2.1 gives $m_\varepsilon-m=O(\varepsilon)$ for fixed $x$.
Thus $H(m_\varepsilon)=H(m)+o(\varepsilon)$, and
\[
B(F_{\varepsilon,x})
=(1-\varepsilon)H(m_\varepsilon)
+\varepsilon|x-m_\varepsilon|
=B(F)+\varepsilon(|x-m|-B(F))+o(\varepsilon).
\]
Therefore $\IF(x;B,F)=|x-m|-B(F)$.
At the Laplace model, $m=\mu$ and $B(F)=b$, giving unbounded
influence $|x-\mu|-b$. In contrast, the location influence is
$b\sgn(x-\mu)$, which is bounded.

\paragraph{Proof of replacement breakdown.}
Replace one observation by $M\to\infty$, keeping an original
observation $x_j$ fixed. For every center $t$,
\[
|M-t|+|x_j-t|\geq|M-x_j|.
\]
Consequently,
\[
\widehat b_M=\frac1n\min_t\sum_i|x_i'-t|
\geq\frac{|M-x_j|}{n}\longrightarrow\infty.
\]
One replacement therefore causes scale explosion.
To include implosion, let $r$ be the maximum multiplicity of an
original value. Replacing $n-r$ observations by that value makes
$\widehat b=0$. With fewer replacements, at least $r+1$ original
observations remain, including two distinct values. If $d>0$ is
the smallest positive separation of distinct original values, then
$\widehat b'\geq d/n$, preventing implosion.
Thus the implosion fraction is $(n-r)/n$, and the explosion fraction
is $1/n$. Their minimum is $1/n$.

\paragraph{Final answer.}
\[
\boxed{\widehat\mu=\operatorname{median}(x_1,\ldots,x_n),\qquad
\widehat b=\frac1n\sum_i|x_i-\widehat\mu|.}
\]
\[
\boxed{\IF(x;B,F_{\mu,b})=|x-\mu|-b,\quad \GES=\infty,
\qquad \epsilon_n^*(\widehat b;x)=\frac1n.}
\]
The asymptotic scale breakdown point is zero. For an all-distinct
sample, the implosion-only fraction is $(n-1)/n$. Rescaling the
parameter to $\sigma=\sqrt2 b$ multiplies its influence function by
$\sqrt2$ and leaves its breakdown point unchanged.
\end{solution}

% ============================================================
\question
Find the MLEs of $\beta$ and $\sigma^2$ for
$Y_i=X_i^\top\beta+\epsilon_i$, with independent
$\epsilon_i\sim N(0,\sigma^2)$, using exponential-family theory.

\begin{solution}
\paragraph{Assumptions.}
Let $X$ be the fixed $n\times p$ design matrix, of full column rank,
with $n>p$. The parameter space is $\R^p\times(0,\infty)$.

\paragraph{Proof.}
The joint Gaussian density can be written in canonical form using
\[
\eta=\frac\beta{\sigma^2},\qquad \tau=-\frac1{2\sigma^2}<0:
\]
\[
f(y;\eta,\tau)
=\exp\{\eta^\top X^\top y+\tau y^\top y-A(\eta,\tau)\},
\]
where
\[
A(\eta,\tau)=\frac n2\log\left(\frac\pi{-\tau}\right)
-\frac{\eta^\top X^\top X\eta}{4\tau}.
\]
The natural space $\R^p\times(-\infty,0)$ is open. This family is
minimal: if $a^\top X^\top y+b\,y^\top y$ is constant for all $y$,
its Hessian forces $b=0$, and full column rank then forces $a=0$.
Thus $A$ is strictly convex and the canonical log-likelihood strictly
concave. An interior solution of the canonical moment equations is
the unique maximizer. The parameter transformation is bijective.

For the canonical statistic $S(Y)=(X^\top Y,Y^\top Y)$, the
exponential-family identity $\nabla A=\E S(Y)$ yields
\[
X^\top y=\E[X^\top Y]=X^\top X\beta,
\]
\[
y^\top y=\E[Y^\top Y]=\beta^\top X^\top X\beta+n\sigma^2.
\]
The first equation gives the least-squares coefficient, and the
second gives the residual sum of squares divided by $n$:
\[
y^\top y-\widehat\beta^\top X^\top X\widehat\beta
=\|y-X\widehat\beta\|^2.
\]
If the residual sum of squares is positive, the solution is interior
and hence is the MLE. If it is zero, an exact fit has likelihood
unbounded as $\sigma^2\downarrow0$.

\paragraph{Final answer.}
For $y\notin\operatorname{col}(X)$,
\[
\boxed{\widehat\beta=(X^\top X)^{-1}X^\top y,\qquad
\widehat\sigma^2=\frac1n\|y-X\widehat\beta\|^2.}
\]
For $y\in\operatorname{col}(X)$, no positive-variance MLE exists.
The MLE uses divisor $n$, not the unbiased-estimation divisor $n-p$.
\end{solution}

% ============================================================
\question
Verify the adaptively weighted least-squares representation of a
regression M-estimate of $\beta$.

\begin{solution}
\paragraph{Assumptions.}
Minimize
\[
L(\beta)=\sum_i\rho(u_i(\beta)),\qquad
u_i(\beta)=\frac{y_i-x_i^\top\beta}{s},\quad s>0,
\]
with scale held fixed. Let $\rho$ be even and differentiable,
$\psi=\rho'$, and suppose $\psi(u)/u$ has a finite continuous
extension at zero. For a weighted least-squares minimization,
assume nonnegative weights and a nonsingular weighted design.

\paragraph{Proof.}
The first-order condition is
\[
\sum_i x_i\psi(u_i(\widehat\beta))=0.
\]
Define
\[
w(u)=\begin{cases}\psi(u)/u,&u\ne0,\\\psi'(0),&u=0,\end{cases}
\qquad
W(\beta)=\diag(w(u_1(\beta)),\ldots,w(u_n(\beta))).
\]
Since $\psi(u)=w(u)u$, the estimating equation becomes
\[
X^\top W(\widehat\beta)(y-X\widehat\beta)=0.
\]
Solving these weighted normal equations gives the desired implicit
representation. Freezing weights at the current iterate gives IRLS.
For Huber and Tukey scores respectively,
\[
w_c^{\mathrm H}(u)=\begin{cases}1,&|u|\leq c,\\c/|u|,&|u|>c,\end{cases}
\qquad
w_c^{\mathrm T}(u)=\left(1-\frac{u^2}{c^2}\right)^2\ind\{|u|\leq c\}.
\]
For these losses, $w$ is even, nonnegative, and nonincreasing in
$|u|$. Let $g(v)=\rho(\sqrt v)$ for $v\geq0$. Since
$g'(v)=w(\sqrt v)/2$ for $v>0$, $g$ is concave, and its tangent
inequality implies
\[
\rho(u)\leq\rho(u_0)+\frac12w(u_0)(u^2-u_0^2).
\]
The extension at $u_0=0$ follows by continuity. Summing these
inequalities produces a quadratic majorizer touching $L$ at the
current iterate. Exact WLS minimization therefore cannot increase
$L$, provided scale is fixed. This does not guarantee convergence
to a global minimizer for a nonconvex loss.

\paragraph{Final answer.}
\[
\boxed{\widehat\beta
=[X^\top W(\widehat\beta)X]^{-1}X^\top W(\widehat\beta)y,}
\]
with iterative update
\[
\boxed{\beta^{(k+1)}
=[X^\top W(\beta^{(k)})X]^{-1}X^\top W(\beta^{(k)})y.}
\]
The weights depend on fitted residuals, so the representation is a
fixed-point identity, not a closed-form solution. Zero weights may
make the design singular. Residual weighting alone does not generally
control high-leverage design points. If scale is estimated separately,
it is held fixed during this argument; alternating scale updates
require a separate scale equation and convergence analysis.
\end{solution}

\end{questions}
\end{document}
